\clearpage\begingroup\SourceLeaf{671b}
\ArithmeticHead{36}{L}{H. GEPPERT, ÜBER GAUSS’ ARBEITEN ZUR MECHANIK UND POTENTIALTHEORIE.}
\makeatletter\tagsleft@true\let\veqno\@@leqno\makeatother
\renewcommand{\thefootnote}{\arabic{footnote})}\setcounter{footnote}{64}
\begin{center}\textls[100]{2. Die \textsc{Gauss}sche Abhandlung.}\end{center}
\noindent\hspace*{\parindent}Dies waren die Arbeiten, die über das Ellipsoidproblem vorlagen, bevor
\textsc{Gauss} es mit neuen Methoden in Angriff nahm. Bei dem etwa 14tägigen
Besuch auf Seeberg bei Gotha, den \textsc{Gauss} im September 1812 \textsc{Lindenau} abstattete,
mag ein Gespräch über den Unterschied im Niveau des Roten und
des Mittelländischen Meeres \textsc{Gauss} den ersten Anstoss zur Bearbeitung des
Attraktionsproblems der Sphäroide gegeben haben\footnote{Vergl. \textsc{Galle}, \textit{Über die geodätischen Arbeiten von Gauss.} Werke XI, 2, S. 43.}. Am 26. Sept. 1812
macht \textsc{Gauss} die Tagebucheintragung\footnote{Werke X, 1, S. 570.}: »Theoriam attractionis sphaeroidis
elliptici in puncta extra solidum sita prorsus novam invenimus«; und gleich
darauf am 15. Okt. 1812: »Etiam partes reliquas eiusdem theoriae per methodum
novam mirae simplicitatis absolvimus«. Von diesen Arbeiten berichtet
\textsc{Gauss} alsbald an seine Freunde \textsc{Gerling}, \textsc{Olbers} und \textsc{Schumacher}\footnote{Vergl. Brief an \textsc{Gerling} vom 15. Nov. 1812, \textsc{Gauss-Gerling}-Briefwechsel, S. 7; Brief an \textsc{Olbers}
vom 8. April 1813, \textsc{Gauss-Olbers}-Briefwechsel I, S. 515; Brief an \textsc{Schumacher} vom 31. Dez. 1812, \textsc{Gauss-Schumacher}-Briefwechsel 1, S. 95.}; letzterer
machte ihn (vergl. die Briefe vom 23. Jan. 1813 und 8. Febr. 1813) auf die
Arbeiten \textsc{Legendres}, bezw. dessen Besprechungen des \textsc{Ivory}schen Aufsatzes aufmerksam.
Am 5. Nov. 1812 berichtet \textsc{Gauss} an \textsc{Laplace} über seine Resultate\footnote{Werke X, 1, S. 378 f.}
und schickt ihm einen Auszug, um ihn dem Institut vorlegen zu lassen;
dieser, heute in der \textsc{Mittag-Leffler}-Bibliothek befindliche Auszug ist Werke
XII, S. 110—114 abgedruckt. \textsc{Gauss} hatte bei seinen Arbeiten wohl die
Untersuchungen von \textsc{Laplace} und seinen Vorgängern berücksichtigt, dagegen
war ihm die Abhandlung von \textsc{Ivory} entgangen; \textsc{Laplace} machte ihn im Briefe
vom 20. Nov. 1812\footnote{Werke X, 1, S. 380—381.} auf diese Arbeit aufmerksam, und kurz darauf, nämlich
\clearpage\SourceLeaf{672}\ArithmeticHead{37}{R}{III. DIE ANZIEHUNG DER ELLIPSOIDE.}
\noindent am 18. März 1813, legte \textsc{Gauss} seinen zusammenfassenden Bericht der Göttinger
Gesellschaft der Wissenschaften vor\footnote{Göttinger Gelehrte Anzeigen vom 5. April 1813; Werke V, S. 279—285.}; im gleichen Jahre erschien die ausführliche
Abhandlung \textit{Theoria attractionis corporum sphaeroidicorum ellipticorum
homogeneorum methodo nova tractata}\footnote{Comment. soc. reg. scient. Gottingensis, 2, 1813; Werke V, S. 1—22.}.

Die \textsc{Gauss}sche Arbeit zerfällt in zwei Teile; der erste und wichtigere
enthält eine Reihe allgemeiner Integralsätze und potentialtheoretischer Theoreme,
die z.\,T. auf früheren funktionentheoretischen Untersuchungen von \textsc{Gauss}
fussen, und deren letzte Ausbeutung er erst später in der Potentialtheorie
unternommen hat; über die Entstehung dieser Sätze werden wir weiter unten
berichten. Der zweite Teil der Abhandlung enthält dann die spezielle Anwendung
auf das Ellipsoidproblem.

Der erste Teil enthält sechs Sätze, die in zwei völlig parallele Satzgruppen
zerfallen. Die drei ersten Sätze beruhen auf dem Gedanken einer prismatischen
Zerlegung des Raumes, während die folgenden drei sich auf eine Aufspaltung
in Elementarkegel gründen. \(O\) bedeute im folgenden eine geschlossene
reguläre Oberfläche, \(B\) den davon eingeschlossenen Körper. Die gedachten
Sätze sind folgende:
\[
\iint\limits_O F(y,z)\cos(n,x)\,dO=0,\tag*{\text{\textls[80]{Satz 1.}}}
\]
\noindent worin \(F\) eine rationale Funktion von \(y,z\) und \(n\) die äussere Normale von \(O\)
bezeichnet.

\textls[80]{Satz 2.}
\begingroup\small
\[
\text{Volumen von }B=\iint\limits_O x\cos(x,n)\,dO=\iint\limits_O y\cos(y,n)\,dO=\iint\limits_O z\cos(z,n)\,dO.
\]
\endgroup
\noindent\hspace*{\parindent}\textls[80]{Satz 3.} Drückt \(f(r)\) das Anziehungsgesetz aus, und ist \(F(r)=\int f(r)\,dr\),
so ist die \(x\)-Komponente der Anziehung des Körpers \(B\) auf einen Punkt
\[
X=\iint\limits_O F(r)\cos(n,x)\,dO.
\]

Die zweite Satzgruppe beginnt mit dem wichtigen »Satze von der scheinbaren
Grösse«.
\clearpage\SourceLeaf{673}\ArithmeticHead{38}{L}{H. GEPPERT, ÜBER GAUSS’ ARBEITEN ZUR MECHANIK UND POTENTIALTHEORIE.}
\textls[80]{Satz 4.} Ist \(M\) der Aufpunkt, \(P\) der variable Punkt auf \(O\), und \(r\) der
Abstand der beiden, so ist
\[
-\iint\limits_O\frac{1}{r^2}\cos(PM,n)\,dO=0,\;2\pi,\;4\pi,
\]
\noindent jenachdem \(M\) ausserhalb, auf oder innerhalb \(O\) liegt, und dieses Integral bedeutet
geometrisch die scheinbare Grösse von \(O\) von \(M\) aus gesehen.
\[
\text{Volumen von }B=-\tfrac13\iint\limits_O r\cos(PM,n)\,dO,\tag*{\text{\textls[80]{Satz 5.}}}
\]
\noindent wobei der Aufpunkt \(M\) beliebig liegen kann.

\textls[80]{Satz 6.} Ist \(\Phi(r)=\int r^2f(r)\,dr\), so ist die \(x\)-Komponente der Anziehung
des Körpers \(B\) auf \(M\)
\[
X=\iint\limits_O\frac{\Phi(r)}{r^2}\cos(PM,n)\cos(PM,x)\,dO
\left\{\begin{aligned}&+\pi\Phi(0)\cos(n,x)\\&+0,\end{aligned}\right.
\]
\noindent jenachdem \(M\) auf \(O\) liegt oder nicht. Zwar ist beim \textsc{Newton}schen Anziehungsgesetz
dieses letzte Zusatzglied unwichtig, aber wichtig erscheint es vom Standpunkte
der Mechanik aus, entspricht es doch einer normal auf \(O\) wirkenden
Kraft, die bei anderen Attraktionsgesetzen sogar unendlich gross werden kann.

Der zweite Teil der \textsc{Gauss}schen Abhandlung gibt die Anwendung der
gefundenen Sätze auf das Ellipsoidproblem. \textsc{Gauss} verwendet hier wie \textsc{Ivory}
eine Parameterdarstellung der Fläche:
\[
x=A\cos p,\quad y=B\sin p\cos q,\quad z=C\sin p\sin q,
\]
\noindent wobei die gleiche Bezeichnung der Parameter wie 1827 in seiner Flächentheorie
auffällt. Setzt man die erste Attraktionskomponente \(X=ABC\xi\), so ergibt
der Satz 3 den Ausdruck
\[
\xi=\iint\limits_O\frac{\sin p\cos p}{Ar}\,dp\,dq.\tag{1}
\]
\noindent Die Integration leistet \textsc{Gauss} durch einen Kunstgriff, der auf der Kenntnis
des \textsc{Mac-Laurin}schen Satzes beruht; er variiert nämlich die Halbachse \(A=\alpha\)
und die beiden andern Halbachsen \(\beta\) und \(\gamma\) so, dass die Ellipsoide konfokal,
also \(\alpha^2-\beta^2\) und \(\alpha^2-\gamma^2\) konstant bleiben. Die Differentiation des Ausdruckes
\clearpage\SourceLeaf{674}\ArithmeticHead{39}{R}{III. DIE ANZIEHUNG DER ELLIPSOIDE.}
\noindent (1) nach \(\alpha\) und die Anwendung des 4. und 6. Integralsatzes ergibt dann
\begingroup\small
\[
\frac{d\xi}{d\alpha}=\left\{\begin{array}{c}\dfrac{-4\pi a}{\alpha^2\beta\gamma}\\0\end{array}\right.,
\quad\text{jenachdem }M\quad
\begin{array}{l}\text{innerhalb}\\\text{ausserhalb}\end{array}
\quad\text{des Ellipsoids liegt.}\tag{2}
\]
\endgroup
\noindent Daraus folgt im zweiten Falle die Konstanz von \(\xi\), d.\,h. der \textsc{Mac-Laurin}sche
Satz, der solange gilt, als \(M\) ausserhalb oder auf dem Ellipsoid liegt. Mithin
reduziert sich die Berechnung von \(X\) auf zwei Aufgaben:

1. Bestimmung eines zum gegebenen konfokalen Ellipsoids durch den
Aufpunkt \(M\),

2. Bestimmung der Anziehung des Ellipsoids auf einen Punkt der Oberfläche.

Der erste Punkt erledigt sich durch Auflösung einer kubischen Gleichung
mit nur einer reellen Wurzel. Der zweite Punkt erfordert die Integration
des ersten Ausdruckes (2), wobei \(\beta\) und \(\gamma\) ebenfalls als Funktionen von \(\alpha\) anzusehen
sind; es resultieren für \(\xi,\eta,\zeta\) elliptische Integrale 2. Gattung, die
sich im Falle des Sphäroids auf elementare Funktionen zurückführen lassen.
Der \textsc{Newton}sche Satz ist sofort aus den Endformeln abzulesen.

In einer handschriftlichen Bemerkung, die Werke V, S. 285—286 wiedergegeben
ist, berechnet \textsc{Gauss} nach den gleichen Methoden das Potential des
Ellipsoids im Aufpunkte \(M\)
\[
V=ABCw=-\tfrac12\iint\limits_O\cos(PM,n)\,dO.
\]
\noindent Die gleiche Variation, die oben angewandt wurde, zeigt dann, dass, falls \(M\)
ausserhalb liegt, \(\frac{dw}{d\alpha}=0\) ist, während für den Fall eines inneren Aufpunktes
sich \(w\) als Summe von vier elliptischen Integralen darstellen lässt.

Was der \textsc{Gauss}schen Lösung des Ellipsoidproblems, verglichen mit der
seiner Vorgänger, eigentümlich ist, ist erstens die grosse Eleganz, die durch
die konsequente Anwendung seiner 6 Integralsätze erreicht wird, und zweitens
die Vermeidung jeglicher Reihenentwicklungen und der damit verbundenen
schwierigen Konvergenzuntersuchungen.
\clearpage\endgroup
